\section{Conclusions and Roadmap}
\label{sec:conclusions}

\subsection{Summary of Theoretical Contributions}

In Part 1 of this series, we addressed the long-standing Richardson threshold paradox by constructing a multiscale, fast-slow geometric framework for the nocturnal stable boundary layer (SBL). By reformulating SBL dynamics using Geometric Singular Perturbation Theory (GSPT) on a desingularized coordinate chart \citep{Fenichel1979, Kuehn2015}, we established the following structural results:

\begin{enumerate}
    \item \textbf{Existence and Smoothness of the Critical Manifold ($\mathcal{S}_0$):} We proved that setting $\epsilon_1 = 0$ yields a 3-dimensional smooth critical manifold embedded in 5D state space $\Omega_0 \subset \mathbb{R}^5$, which decomposes into a laminar extinction sheet $\mathcal{S}_0^0$ and a curved turbulent sheet $\mathcal{S}_0^+$.

    \item \textbf{Fold Characterization and Fast Degeneracy:} Loss of normal hyperbolicity along $\mathcal{S}_0^+$ is governed by the fast Jacobian determinant $\det J_{\mathrm{fast}} = 0$, defining a 2-dimensional invariant fold surface $\mathcal{C}_{\mathrm{fold}} \subset \Omega_0$ where turbulent fast adjustment collapses \citep{VanDeWiel2002}.

    \item \textbf{Resolution of the Richardson Threshold Paradox:} We proved the Fold-Surface Projection Theorem (Theorem~3). Under incomplete diagnostic observation operators $\boldsymbol{\Pi}_{\mathrm{obs}}$ that omit soil thermal states, the invariant 2D fold surface $\mathcal{C}_{\mathrm{fold}}$ maps to a 1-dimensional threshold curve $\Gamma_{\mathrm{fold}}$ in diagnostic space $(S, Ri)$. Field campaigns (e.g., CASES-99 and SHEBA) trace distinct 1D site profile curves $\gamma_{\mathrm{site}} = \mathcal{S}_0^+ \cap \Sigma_{\mathrm{site}}$ that intersect $\Gamma_{\mathrm{fold}}$ at different scalar values \citep{Poulos2002CASES99, Uttal2002SHEBA}. Thus, critical Richardson numbers are not universal physical constants, but site-specific scalar projections of an invariant higher-dimensional geometry.
\end{enumerate}

\subsection{Roadmap for Parts 2 and 3}

The invariant geometric skeleton developed in Part~1 provides the mathematical foundation for operational parameterizations and data-assimilative validation in subsequent papers:

\begin{itemize}
    \item \textbf{Part 2: Empirical Validation, Parameter Estimation, and Campaign Phase Space Mapping.}
    Part~2 projects high-frequency turbulence data from CASES-99, SHEBA, and Cabauw directly onto the desingularized coordinate chart $\Omega_0$. Using modern manifold learning and Bayesian parameter estimation, we reconstruct empirical critical manifolds $\mathcal{S}_0^{\mathrm{obs}}$, verify the location of $\mathcal{C}_{\mathrm{fold}}$, and quantify site-specific constraint manifolds $\Sigma_{\mathrm{site}}$.

    \item \textbf{Part 3: Dynamic Threshold Closures and WRF Implementation.}
    Part~3 translates the analytical fold formulas derived here ($H_{\max}$ and $Ri_{\mathrm{fold}}(T_s, T_g)$) into a non-local boundary-layer scheme for numerical weather prediction (NWP) \citep{Holtslag2013}. We demonstrate that incorporating fold-aware dynamic thresholds eliminates systemic runaway cooling biases in single-column models (SCM) and regional weather research forecast models (WRF).
\end{itemize}